\subsection{Sets, Maps, Relations, and Universal Properties}

The elementary set, order, choice, function, and bijection vocabulary used here is collected in the foundational prerequisites. This subsection continues with set algebra, universal properties, and equivalence-based structure.

\begin{definition} \label{def:equivalence-relation} \label{def:equivalence-classes-invariants-canonical-forms} \hlt{Equivalence Relations, Classes, Invariants, and Canonical Forms}\\
Let $\sim$ be a relation on a set $X$.
\begin{enumerate}[label=(\roman*)]
\setlength{\itemsep}{0pt}
\item $\sim$ is an \hlt{equivalence relation} if, for all $x,y,z\in X$, $x\sim x$, $x\sim y\Rightarrow y\sim x$, and $(x\sim y\text{ and }y\sim z)\Rightarrow x\sim z$.
\item If $\sim$ is an equivalence relation, the \hlt{equivalence class} of $x\in X$ is $[x]\defeq\{y\in X:y\sim x\}$.
\item A function $f:X\to Y$ is an \hlt{invariant} of $\sim$ if $x\sim y\Rightarrow f(x)=f(y)$, and a \hlt{complete invariant} if $x\sim y\Longleftrightarrow f(x)=f(y)$.
\item A subset $\mathcal{C}\subseteq X$ is an \hlt{set of canonical forms} for $\sim$ if each equivalence class meets $\mathcal{C}$ in exactly one element.
\end{enumerate}
\end{definition}
